%-------------------------------------------------------------------------------
% Targeted Resume -- Chenyu Gu
% Role: Research Scientist Intern, AI-Driven Neural Interfaces (PhD)
%       Meta Reality Labs (CTRL-labs) -- Burlingame, CA
%
% Designed to fit on ONE page. Build: pdflatex resume-meta-neural-interfaces.tex
%-------------------------------------------------------------------------------

\documentclass[letterpaper,10pt]{article}

%----------------------------------------------------
%                     Packages
%----------------------------------------------------
\usepackage{latexsym}
\usepackage[empty]{fullpage}
\usepackage{titlesec}
\usepackage{marvosym}
\usepackage[usenames,dvipsnames]{color}
\usepackage{enumitem}
\usepackage[hidelinks]{hyperref}
\usepackage{fancyhdr}
\usepackage[english]{babel}
\usepackage{tabularx}
\usepackage{hyphenat}
\usepackage{fontawesome}
\input{glyphtounicode}

% Ensure the generated PDF is machine readable / ATS parsable
\pdfgentounicode=1

%----------------------------------------------------
%                    Page Setup
%----------------------------------------------------
\pagestyle{fancy}
\fancyhf{}
\fancyfoot{}
\renewcommand{\headrulewidth}{0pt}
\renewcommand{\footrulewidth}{0pt}

\addtolength{\oddsidemargin}{-0.62in}
\addtolength{\evensidemargin}{-0.62in}
\addtolength{\textwidth}{1.24in}
\addtolength{\topmargin}{-0.75in}
\addtolength{\textheight}{1.5in}

\urlstyle{same}
\raggedbottom
\raggedright
\setlength{\tabcolsep}{0in}

% Section formatting
\titleformat{\section}{
  \vspace{-7pt}\scshape\raggedright\large
}{}{0em}{}[\color{black}\titlerule \vspace{-6pt}]

%----------------------------------------------------
%                  Custom Commands
%----------------------------------------------------
\newcommand{\resumeItem}[1]{
  \item\small{{#1 \vspace{-2.5pt}}}
}

\newcommand{\resumeSubheading}[4]{
  \vspace{-1pt}\item
  \begin{tabular*}{0.98\textwidth}[t]{l@{\extracolsep{\fill}}r}
    \textbf{#1} & \small#2 \\
    \textit{\small#3} & \textit{\small #4} \\
  \end{tabular*}\vspace{-6pt}
}

\newcommand{\resumeEduHeading}[4]{
  \vspace{-1pt}\item
  \begin{tabular*}{0.98\textwidth}[t]{l@{\extracolsep{\fill}}r}
    \textbf{#1} & \small#2 \\
    \textit{\small#3} & \textit{\small #4} \\
  \end{tabular*}\vspace{-6pt}
}

\renewcommand\labelitemii{$\vcenter{\hbox{\tiny$\bullet$}}$}

\newcommand{\resumeSubHeadingListStart}{\begin{itemize}[leftmargin=0.15in, label={}]}
\newcommand{\resumeSubHeadingListEnd}{\end{itemize}}
\newcommand{\resumeItemListStart}{\begin{itemize}[leftmargin=0.28in, topsep=1pt, itemsep=0pt, parsep=0pt]}
\newcommand{\resumeItemListEnd}{\end{itemize}\vspace{-4pt}}

%-------------------------------------------------------------------------------
\begin{document}
%-------------------------------------------------------------------------------

%----------------------------------------------------
%                      Header
%----------------------------------------------------
\begin{center}
  \textbf{\Large \scshape Chenyu Gu} \\ \vspace{2pt}
  \footnotesize
  \faMobile \hspace{.5pt} \href{tel:14046374484}{+1 (404) 637 4484}
  $|$
  \faAt \hspace{.5pt} \href{mailto:gu.c@wustl.edu}{gu.c@wustl.edu}
  $|$
  \faLinkedinSquare \hspace{.5pt} \href{https://www.linkedin.com/in/chenyu-gu-38337226b/}{LinkedIn}
  $|$
  \faMapMarker \hspace{.5pt} Saint Louis, MO
\end{center}
\vspace{-10pt}

%----------------------------------------------------
%                  Research Focus
%----------------------------------------------------
\section{Research Focus}
\vspace{1pt}
{\small Biomedical Engineering Ph.D.\ student working on \textbf{neural interfaces and wearable biosignal
decoding}: turning noisy, multi-channel body-worn sensor data (\textbf{EMG}, \textbf{IMU},
stimulation-evoked neuromotor responses) into models of user intent --- hardware and real-time firmware
through signal processing, deep learning in \textbf{PyTorch}, and human-subject validation.
Seeking a 12--24 week research internship; I intend to return to my Ph.D.\ program upon completion.}

%----------------------------------------------------
%                   Education
%----------------------------------------------------
\vspace{2pt}
\section{Education}
\vspace{2pt}
\resumeSubHeadingListStart

\resumeEduHeading
{Washington University in St.\ Louis}{Saint Louis, MO}
{\textbf{Ph.D.} in Biomedical Engineering --- Seanez Lab (neural interfaces); \textit{Advisor:} Dr.\ Ismael Seanez}{Aug 2026 -- Present}

\resumeEduHeading
{Georgia Institute of Technology}{Atlanta, GA}
{\textbf{M.S.} in Bioengineering (ECE track); \textbf{GPA 4.00/4.00}; \textit{Advisor:} Dr.\ Jaydev P.\ Desai}{Aug 2024 -- May 2026}

\resumeEduHeading
{Southern University of Science and Technology}{Shenzhen, China}
{\textbf{B.S.} in Intelligent Medical Engineering; \textbf{GPA 3.82/4.00}; \textit{Advisor:} Dr.\ Mingming Zhang}{Aug 2020 -- Jun 2024}

\vspace{1pt}
\resumeItem{\footnotesize \textit{Coursework:} Machine Learning, Digital Signal Processing, Signals and
Systems, Neural Engineering and Brain--Computer Interfaces, Digital Control, Robotics. Exchange semester
at KTH Royal Institute of Technology.}

\resumeSubHeadingListEnd

%----------------------------------------------------
%                Research Experience
%----------------------------------------------------
\vspace{-3pt}
\section{Research Experience}
\vspace{2pt}
\resumeSubHeadingListStart

%---------- Seanez Lab ----------
\resumeSubheading
{Graduate Research Assistant --- Seanez Lab, Washington University}{Saint Louis, MO}
{Multi-channel EMG decoding and spatiotemporal neurostimulation}{Aug 2026 -- Present}
\resumeItemListStart
\resumeItem{Lead a 16-participant human study on \textbf{multielectrode transcutaneous spinal cord
stimulation arrays}, optimizing spatial electrode patterns to steer targeted joint torque --- multi-channel
array processing over evoked neuromotor signals.}
\resumeItem{Built Python pipelines for \textbf{16-channel wireless surface EMG}: filtering, stimulus
artifact handling, stimulus-locked epoching, and quantification of posterior root--muscle reflexes via
post-activation depression.}
\resumeItem{Developed \textbf{real-time} frequency-modulated stimulation and monitoring software tracking
evoked muscle responses alongside autonomic nociception and pain profiles.}
\resumeItem{Time-synchronized high-density EMG, dynamometer torque, and 3D kinematics into one
\textbf{multimodal dataset}; modeled recruitment selectivity with linear mixed-effects models over a
15-session crossover design.}
\resumeItemListEnd

%---------- Motion capture / deep learning ----------
\resumeSubheading
{Project Leader --- Brain-Robot Rehabilitation Lab, SUSTech}{Shenzhen, China}
{Deep learning for sparse wearable motion capture (IMU + UWB)}{Apr 2022 -- May 2024}
\resumeItemListStart
\resumeItem{Built a hybrid \textbf{IMU + UWB} wearable capture system and trained \textbf{deep learning
models reconstructing full-body kinematics from sparse sensor sets} --- sequence modeling of noisy,
multi-channel time-series under real wearing conditions.}
\resumeItem{Owned the full stack --- Python acquisition and streaming, model training and evaluation,
real-time Unity visualization --- and validated it on human subjects during rehabilitation exercises.}
\resumeItem{First author of a systematic review on IMU motion capture (\textbf{Annual Best Paper Award});
co-inventor on an issued patent for the motion capture method.}
\resumeItemListEnd

%---------- Wearable inertial navigation ----------
\resumeSubheading
{Project Leader --- Brain-Robot Rehabilitation Lab, SUSTech}{Shenzhen, China}
{Wearable 9-axis IMU sensing system, hardware to algorithm}{Feb 2023 -- Jun 2024}
\resumeItemListStart
\resumeItem{Co-designed circuit, embedded C firmware, and host software for a body-worn inertial tracking
system with real-time orientation feedback.}
\resumeItem{Validated accuracy and repeatability against optical tracking and a robot arm, then in the
operating room; published in \textit{IEEE Trans.\ Instrumentation and Measurement}.}
\resumeItemListEnd

%---------- KTH ----------
\resumeSubheading
{Visiting Researcher --- Intelligent Heart Technology Lab, KTH}{Stockholm, Sweden}
{Sensor fusion for gait and physiological signal analysis}{Oct 2023 -- Jan 2024}
\resumeItemListStart
\resumeItem{Implemented \textbf{Kalman filtering for IMU orientation estimation} and sensor calibration;
developed gait-parameter and center-of-mass extraction algorithms.}
\resumeItem{Related the extracted movement features to cardiovascular measures, linking wearable kinematic
signals to physiological state.}
\resumeItemListEnd

%---------- RoboMed ----------
\resumeSubheading
{Graduate Researcher --- RoboMed Lab, Georgia Tech}{Atlanta, GA}
{Data-driven modeling and control of a continuum robot}{Oct 2024 -- May 2026}
\resumeItemListStart
\resumeItem{Modeled a polymer hydraulic continuum robot from first principles, refined it with a
\textbf{data-driven hysteresis model}, and implemented open- and closed-loop control.
\textbf{ISMR 2026 Best Student Paper Award}.}
\resumeItemListEnd

\resumeSubHeadingListEnd

%----------------------------------------------------
%                 Technical Skills
%----------------------------------------------------
\vspace{-4pt}
\section{Technical Skills \& Teaching}
\vspace{2pt}
\resumeItemListStart
\resumeItem{\textbf{Languages \& Tools:} Python (NumPy, SciPy, pandas, scikit-learn, \textbf{PyTorch}),
C/C++ (embedded, real-time), MATLAB/Simulink, Rust, Java; Git, Linux, Docker, Jupyter.}
\resumeItem{\textbf{Signals \& ML:} Multi-channel biosignal processing (EMG, IMU, ECG), filtering and
artifact rejection, Kalman filtering and sensor fusion, time-series/sequence modeling, deep learning,
dimensionality reduction and clustering, mixed-effects models and experimental design.}
\resumeItem{\textbf{Systems \& Hardware:} Wearable sensor system design, PCB and analog circuits,
microcontroller firmware, real-time streaming and multi-device synchronization, Unity visualization,
optical/inertial motion capture, isokinetic dynamometry.}
\resumeItem{\textbf{Teaching:} Graduate TA, \textbf{Machine Learning (CS 7641)}, Georgia Tech --- office
hours, review sessions, and project grading for 300+ graduate students (Aug 2025 -- May 2026).}
\resumeItem{\textbf{Research output:} 4 peer-reviewed papers (2 best-paper awards) and 2 issued patents;
full list available in my full CV.}
\resumeItemListEnd

\end{document}
